\section{Solved Problems on Advanced Interval Estimation and Sample Size}

\subsection{Problem Statements}

\subsubsection*{1. Fundamental Level (Simple / Direct)}

\begin{problema}
	\label{prob:n_media}
	A research team wishes to estimate the average time a user spends interacting with a virtual educational platform before logging out. From preliminary data on similar platforms, it is known that the standard deviation of session durations is roughly $\sigma = 18$ minutes.
	\begin{enumerate}
		\item What sample size is required to estimate the population mean $\mu$ with a maximum tolerable margin of error $E = 2$ minutes at a $95\%$ confidence level?
		\item If the budget only allows monitoring a maximum of $n = 150$ sessions, what would be the actual margin of error achieved at $95\%$ confidence?
	\end{enumerate}
\end{problema}

\begin{problema}
	\label{prob:ic_varianza}
	A financial analytics team monitors the daily volatility of a technology stock market index. Over a trading month of $n = 21$ business days, logarithmic returns exhibit a sample standard deviation $S = 1.45\%$. Assuming that daily returns follow a normal distribution, construct a $90\%$ confidence interval for the true population variance $\sigma^2$ and for the standard deviation $\sigma$.
\end{problema}

\begin{problema}
	\label{prob:dif_medias_z_avanzado}
	In a data center performance audit, query processing response times are compared between two independent server clusters. From a random sample of $n_1 = 50$ queries on Cluster 1, an average of $\bar{X}_1 = 120.0$ ms with a sample standard deviation $S_1 = 15.0$ ms was obtained. On Cluster 2, a sample of $n_2 = 60$ queries yielded $\bar{X}_2 = 112.0$ ms with $S_2 = 12.0$ ms.
	
	Given that the sample sizes are large ($n_1, n_2 \ge 30$), construct a $95\%$ parametric confidence interval for the difference in average response times $\mu_1 - \mu_2$ using the standard normal distribution $Z$.
\end{problema}

\subsubsection*{2. Operational Level (Intermediate / Developmental)}

\begin{problema}
	\label{prob:dif_medias_indep_iguales}
	A data science engineer evaluates the latency time (in milliseconds) of two cloud database architectures (Architecture A and Architecture B) under similar workloads. They draw two independent random samples with the following results:
	\begin{itemize}
		\item \textbf{Architecture A:} $n_1 = 15$, sample mean $\bar{X}_1 = 48.5$ ms, standard deviation $S_1 = 6.2$ ms.
		\item \textbf{Architecture B:} $n_2 = 12$, sample mean $\bar{X}_2 = 53.8$ ms, standard deviation $S_2 = 5.8$ ms.
	\end{itemize}
	Assuming that the populations are normal and have equal variances ($\sigma_1^2 = \sigma_2^2$), construct a $95\%$ confidence interval for the true difference in mean latencies $\mu_1 - \mu_2$.
\end{problema}

\begin{problema}
	\label{prob:dif_proporciones_ab}
	In an A/B test for an e-commerce website, two versions of a checkout form are randomly shown to visitors. Out of $n_1 = 800$ users exposed to Version A (classic design), $X_1 = 144$ completed a purchase. Out of $n_2 = 850$ users exposed to Version B (single-step minimalist design), $X_2 = 187$ completed a purchase.
	\begin{enumerate}
		\item Construct a $95\%$ confidence interval for the true difference between population conversion rates $p_1 - p_2$.
		\item Do the data justify the decision to permanently adopt Version B?
	\end{enumerate}
\end{problema}

\begin{problema}
	\label{prob:n_proporcion}
	An analytical firm plans to conduct a national survey to determine the proportion $p$ of urban households using generative artificial intelligence services in their daily lives.
	\begin{enumerate}
		\item Calculate the minimum required sample size so that the estimate has a margin of error no greater than $3\%$ ($E = 0.03$) at a $99\%$ confidence level, assuming no prior survey is available (most conservative scenario).
		\item Suppose a previous pilot study suggested that this proportion is around $20\%$ ($\hat{p}^* = 0.20$). How many surveys would be required under this new scenario with the same level of precision and confidence? What percentage savings in sample size does having this prior information represent?
	\end{enumerate}
\end{problema}

\subsubsection*{3. Analytical Level (Advanced / Connection)}

\begin{problema}
	\label{prob:dif_medias_welch}
	In a study on the performance of two optimization algorithms (Algorithm 1 and Algorithm 2) applied to convolutional neural networks, convergence time was measured in seconds. Since the variability of one seems much larger, equal variances are not assumed ($\sigma_1^2 \neq \sigma_2^2$). The obtained data are:
	\begin{itemize}
		\item \textbf{Algorithm 1:} $n_1 = 10$, $\bar{X}_1 = 120.4$ s, $S_1 = 18.5$ s.
		\item \textbf{Algorithm 2:} $n_2 = 14$, $\bar{X}_2 = 105.2$ s, $S_2 = 7.3$ s.
	\end{itemize}
	Construct a $99\%$ confidence interval for $\mu_1 - \mu_2$ using the Welch-Satterthwaite approximation.
\end{problema}

\begin{problema}
	\label{prob:ic_razon_varianzas}
	To verify whether two industrial IoT temperature sensors manufactured by different suppliers have exactly the same variance and precision in their measurements, two independent samples are taken:
	\begin{itemize}
		\item \textbf{Supplier 1:} $n_1 = 16$ measurements, with sample variance $S_1^2 = 0.36$ $^\circ$C$^2$.
		\item \textbf{Supplier 2:} $n_2 = 13$ measurements, with sample variance $S_2^2 = 0.12$ $^\circ$C$^2$.
	\end{itemize}
	Construct a $98\%$ confidence interval for the population variance ratio $\dfrac{\sigma_1^2}{\sigma_2^2}$. Is there statistical evidence to claim that one supplier has higher variability than the other?
\end{problema}

\subsubsection*{4. Challenging Level (Experts / Deep Deduction)}

\begin{problema}
	\label{prob:ic-desaf-fisher}
	\textbf{Derivation of the exact asymptotic interval for Pearson's correlation using Fisher's transformation:}
	The sample linear correlation coefficient $r$ obtained from a bivariate normal sample $(X_i, Y_i)$ of size $n$ has a distribution on the interval $[-1, 1]$ that becomes highly skewed as the true parameter $\rho \to \pm 1$.
	\begin{enumerate}
		\item Theoretically explain why the symmetric standard error technique $r \pm Z_{\alpha/2}\text{SE}(r)$ produces mathematical absurdities or poor coverage when $|\rho|$ is large in small or moderate samples.
		\item Fisher's hyperbolic variance-stabilizing transformation is defined by:
		\begin{align}
			Z(r) = \text{artanh}(r) = \frac{1}{2} \ln\left( \frac{1 + r}{1 - r} \right).
		\end{align}
		Algebraically demonstrate how to invert the symmetric asymptotic confidence interval for $Z(\rho)$ in the transformed scale ($Z(r) \pm Z_{\alpha/2}\frac{1}{\sqrt{n-3}}$) using the hyperbolic tangent function $\tanh(z) = \frac{e^{2z}-1}{e^{2z}+1}$ to obtain explicit limits $[L_\rho, \, U_\rho]$ in the original scale of the correlation coefficient.
		\item In a bioinformatics analysis on collinear gene expression across $n = 28$ patients, a sample correlation coefficient $r = 0.82$ was observed. Exactly compute the $95\%$ confidence interval for $\rho$.
	\end{enumerate}
\end{problema}

\begin{problema}
	\label{prob:ic-desaf-delta-method}
	\textbf{The Delta Method for interval inference on ratios of means:}
	In quantitative economics and data science, one frequently wishes to estimate and construct confidence intervals for ratios of population parameters $\theta = \dfrac{\mu_Y}{\mu_X}$ (for example, average elasticity, cost efficiency, or return on cost), from independent pairs of observations $(X_1, Y_1), \dots, (X_n, Y_n)$ where $\mu_X \neq 0$. The natural ratio estimator is defined as:
	\begin{align}
		\hat{\theta} = \frac{\bar{Y}}{\bar{X}}.
	\end{align}
	\begin{enumerate}
		\item Using a first-order multivariate Taylor series expansion (or multivariate Delta Method) around the exact point $(\mu_X, \mu_Y)$, rigorously prove that the asymptotic variance of $\hat{\theta}$ is given by the expression:
		\begin{align}
			\Var(\hat{\theta}) \approx \frac{1}{n \mu_X^2} \left[ \sigma_Y^2 - 2\theta \sigma_{XY} + \theta^2 \sigma_X^2 \right],
		\end{align}
		where $\sigma_{XY} = \cov(X, Y)$ is the population covariance.
		\item Based on the derived result, systematically state how to construct the $(1-\alpha)$ asymptotic confidence interval for $\theta = \dfrac{\mu_Y}{\mu_X}$ by replacing unknown parameters with consistent sample moments ($\bar{X}, \bar{Y}, s_X^2, s_Y^2, s_{XY}$).
	\end{enumerate}
\end{problema}


\subsection{Detailed Solutions}

\subsubsection*{1. Fundamental Level (Simple / Direct)}

\begin{solproblema}
	\textbf{Part 1: Required sample size ($n$).}
	For a $95\%$ confidence level, the standardized normal value is $Z_{0.025} = 1.96$. The known dispersion is $\sigma = 18$ min and the tolerable margin of error is $E = 2$ min.
	
	Substituting directly into the elementary sample size formula:
	\begin{align}
		n = \left( \frac{Z_{\alpha/2} \cdot \sigma}{E} \right)^2 = \left( \frac{1.96 \times 18}{2} \right)^2 = (1.96 \times 9)^2 = (17.64)^2 = 311.1696.
	\end{align}
	Since sample sizes correspond to discrete counts of entire sessions and the margin of error must strictly not exceed 2 minutes, we apply the ceiling rule by always rounding up:
	\begin{align}
		n \geq \lceil 311.17 \rceil = 312 \text{ sessions.}
	\end{align}
	
	\textbf{Part 2: Margin of error under a fixed budget of $n = 150$.}
	Solving for $E$ from the confidence interval half-width formula for the mean:
	\begin{align}
		E = Z_{\alpha/2} \frac{\sigma}{\sqrt{n}} = 1.96 \times \frac{18}{\sqrt{150}} = \frac{35.28}{12.2474} \approx 2.8806 \text{ minutes.}
	\end{align}
	Consequently, limiting data collection to $150$ sessions decreases the study's precision, yielding an increased margin of error of $2.88$ minutes at $95\%$ confidence.
\end{solproblema}

\begin{solproblema}
	\textbf{Step 1: Degrees of freedom and chi-square quantiles.}
	The degrees of freedom of the sample variance statistic are $\nu = n - 1 = 21 - 1 = 20$.
	For a $90\%$ confidence level ($\alpha = 0.10$), the upper and lower tails exclude an area of $\alpha/2 = 0.05$ and $1-\alpha/2 = 0.95$. Consulting the $\chi^2$ distribution tables with $\nu = 20$:
	\begin{align}
		\chi^2_{0.05, \, 20} &= 31.410, \\
		\chi^2_{0.95, \, 20} &= 10.851.
	\end{align}
	
	\textbf{Step 2: Interval for the variance $\sigma^2$.}
	The numerator of the sum of squares $(n-1)S^2$ is:
	\begin{align}
		(20)(1.45)^2 = 20 \times 2.1025 = 42.05 \text{ (\%}^2\text{).}
	\end{align}
	Substituting into the exact limits of the $\chi^2$ distribution:
	\begin{align}
		\frac{42.05}{31.410} \le \sigma^2 \le \frac{42.05}{10.851} \implies [1.3387, \, 3.8752] \text{ (\%}^2\text{).}
	\end{align}
	
	\textbf{Step 3: Interval for the standard deviation $\sigma$.}
	Since the square root is a strictly increasing function for positive reals, we take the square root across all terms:
	\begin{align}
		\sqrt{1.3387} \le \sigma \le \sqrt{3.8752} \implies [1.157, \, 1.969] \%.
	\end{align}
	
	\textbf{Interpretation:} With $90\%$ confidence, the true daily volatility of the stock index lies between $1.16\%$ and $1.97\%$.
\end{solproblema}

\begin{solproblema}
	The difference between sample average latencies is:
	\begin{align}
		\hat{D} = \bar{X}_1 - \bar{X}_2 = 120.0 - 112.0 = 8.0 \text{ ms.}
	\end{align}
	Because both samples are large ($n_1 = 50, n_2 = 60$), the standard error of the difference is calculated and standardized without assuming equal variances via the Central Limit Theorem:
	\begin{align}
		\text{SE}(\bar{X}_1 - \bar{X}_2) = \sqrt{\frac{S_1^2}{n_1} + \frac{S_2^2}{n_2}} = \sqrt{\frac{15.0^2}{50} + \frac{12.0^2}{60}} = \sqrt{\frac{225.0}{50} + \frac{144.0}{60}} = \sqrt{4.5 + 2.4} = \sqrt{6.9} \approx 2.6268 \text{ ms.}
	\end{align}
	For a $95\%$ confidence level ($Z_{0.025} = 1.96$), the margin of error is:
	\begin{align}
		E = 1.96 \times 2.6268 \approx 5.1485 \text{ ms.}
	\end{align}
	The 95\% parametric confidence interval results in:
	\begin{align}
		8.0 - 5.1485 \le \mu_1 - \mu_2 \le 8.0 + 5.1485 \implies [2.85, \, 13.15] \text{ ms.}
	\end{align}
	\textbf{Interpretation:} With 95\% confidence, queries on Cluster 1 take on average between $2.85$ and $13.15$ milliseconds longer than on Cluster 2, demonstrating in a statistically robust manner that Cluster 2 is faster.
\end{solproblema}

\subsubsection*{2. Operational Level (Intermediate / Developmental)}

\begin{solproblema}
	\textbf{Step 1: Calculation of the pooled sample variance ($S_p^2$).}
	Assuming normality and homoscedasticity ($\sigma_1^2 = \sigma_2^2$), we compute the weighted pooled variance:
	\begin{align}
		S_p^2 &= \frac{(n_1 - 1)S_1^2 + (n_2 - 1)S_2^2}{n_1 + n_2 - 2} = \frac{(15 - 1)(6.2)^2 + (12 - 1)(5.8)^2}{15 + 12 - 2} \\
		&= \frac{14(38.44) + 11(33.64)}{25} = \frac{538.16 + 370.04}{25} = \frac{908.2}{25} = 36.328 \text{ ms}^2.
	\end{align}
	Therefore, the pooled standard deviation is $S_p = \sqrt{36.328} \approx 6.027$ ms.
	
	\textbf{Step 2: Critical quantile from the $t$-distribution with $\nu = 25$.}
	Total degrees of freedom are $\nu = 15 + 12 - 2 = 25$. For 95\% confidence ($\alpha = 0.05$), the table value is:
	\begin{align}
		t_{0.025, \, 25} = 2.060.
	\end{align}
	
	\textbf{Step 3: Interval construction.}
	The difference in means is $\bar{X}_1 - \bar{X}_2 = 48.5 - 53.8 = -5.3$ ms. The margin of error $E$ is calculated as:
	\begin{align}
		E = t_{0.025, \, 25} \cdot S_p \sqrt{\frac{1}{n_1} + \frac{1}{n_2}} = 2.060(6.027) \sqrt{\frac{1}{15} + \frac{1}{12}} = (12.4156) \sqrt{0.1500} = (12.4156)(0.3873) \approx 4.8086 \text{ ms.}
	\end{align}
	The $95\%$ confidence interval for the latency difference is therefore:
	\begin{align}
		-5.3 - 4.809 \le \mu_1 - \mu_2 \le -5.3 + 4.809 \implies [-10.11, \, -0.49] \text{ ms.}
	\end{align}
	\textbf{Interpretation:} Since the entire interval is contained within the negative region ($[-10.11, -0.49]$) and does not touch zero, it is confirmed with 95\% confidence that Architecture A has a significantly lower mean latency (is faster) than Architecture B.
\end{solproblema}

\begin{solproblema}
	\textbf{Part 1: Proportions and standard error.}
	The empirical purchase (conversion) proportions are:
	\begin{align}
		\hat{p}_1 &= \frac{144}{800} = 0.18 \quad (18\%), \\
		\hat{p}_2 &= \frac{187}{850} = 0.22 \quad (22\%).
	\end{align}
	The point difference in conversion rates is $\hat{p}_1 - \hat{p}_2 = 0.18 - 0.22 = -0.04$. We verify that all observed and expected counts of successes and failures ($144, 656, 187, 663$) exceed $10$, confirming that the asymptotic normal approximation is entirely valid.
	
	The standard error of the difference is computed as:
	\begin{align}
		\text{SE}(\hat{p}_1 - \hat{p}_2) &= \sqrt{\frac{\hat{p}_1(1-\hat{p}_1)}{n_1} + \frac{\hat{p}_2(1-\hat{p}_2)}{n_2}} = \sqrt{\frac{0.18(0.82)}{800} + \frac{0.22(0.78)}{850}} \\
		&= \sqrt{\frac{0.1476}{800} + \frac{0.1716}{850}} = \sqrt{0.0001845 + 0.0002019} = \sqrt{0.0003864} \approx 0.01966.
	\end{align}
	For $95\%$ confidence ($Z_{0.025} = 1.96$), the margin of error $E$ is:
	\begin{align}
		E = 1.96 \times 0.01966 \approx 0.0385 \text{ (or 3.85\%).}
	\end{align}
	The $95\%$ confidence interval for $p_1 - p_2$ is:
	\begin{align}
		-0.04 - 0.0385 \le p_1 - p_2 \le -0.04 + 0.0385 \implies [-0.0785, \, -0.0015].
	\end{align}
	
	\textbf{Part 2: Decision on the A/B design.}
	The interval for $p_1 - p_2$ spans from $-7.85\%$ to $-0.15\%$. Since the upper limit is strictly below zero, there is 95\% statistical certainty that Version B's conversion rate ($p_2$) genuinely exceeds Version A's ($p_1$). Therefore, the data fully justify the immediate adoption of Version B.
\end{solproblema}

\begin{solproblema}
	\textbf{Part 1: Conservative scenario ($\hat{p}^* = 0.5$).}
	For 99\% confidence, the normal quantile is $Z_{0.005} = 2.576$. The allowable margin of error is $E = 0.03$. Lacking prior surveys, we maximize the variance using $p^*(1-p^*) = 0.5(0.5) = 0.25$:
	\begin{align}
		n_{\text{max}} = \frac{Z_{\alpha/2}^2}{4 E^2} = \frac{(2.576)^2}{4 (0.03)^2} = \frac{6.635776}{4 (0.0009)} = \frac{6.635776}{0.0036} = 1,843.27.
	\end{align}
	Applying rounding up by integer count: $n = \lceil 1,843.27 \rceil = 1,844$ households in total.
	
	\textbf{Part 2: Scenario with prior information ($\hat{p}^* = 0.20$).}
	Utilizing the pilot study, the estimated variance decreases to $p^*(1-p^*) = 0.20(0.80) = 0.16$:
	\begin{align}
		n = \left( \frac{Z_{\alpha/2}}{E} \right)^2 p^*(1-p^*) = \left( \frac{2.576}{0.03} \right)^2 (0.16) = (85.8667)^2 (0.16) = (7,373.084)(0.16) = 1,179.69.
	\end{align}
	Rounding up to the next integer: $n = \lceil 1,179.69 \rceil = 1,180$ households in the survey.
	
	\textbf{Evaluation of percentage savings:}
	The absolute savings in surveyed households by utilizing prior data science information is $1,844 - 1,180 = 664$ households. The percentage reduction in data collection costs is:
	\begin{align}
		\text{Savings} = \frac{664}{1,844} \times 100\% \approx 36.01\%.
	\end{align}
\end{solproblema}

\subsubsection*{3. Analytical Level (Advanced / Connection)}

\begin{solproblema}
	\textbf{Step 1: Welch-Satterthwaite degrees of freedom ($\nu$).}
	We calculate the sample variances divided by their respective sample sizes:
	\begin{align}
		\frac{S_1^2}{n_1} = \frac{18.5^2}{10} = \frac{342.25}{10} = 34.225, \quad \frac{S_2^2}{n_2} = \frac{7.3^2}{14} = \frac{53.29}{14} \approx 3.8064.
	\end{align}
	Substituting into the analytical approximation for heterogeneous variances:
	\begin{align}
		\nu &= \frac{\left( \dfrac{S_1^2}{n_1} + \dfrac{S_2^2}{n_2} \right)^2}{\dfrac{(S_1^2 / n_1)^2}{n_1 - 1} + \dfrac{(S_2^2 / n_2)^2}{n_2 - 1}} = \frac{(34.225 + 3.8064)^2}{\dfrac{34.225^2}{10 - 1} + \dfrac{3.8064^2}{14 - 1}} = \frac{(38.0314)^2}{\dfrac{1,171.3506}{9} + \dfrac{14.4887}{13}} \\
		&= \frac{1,446.3874}{130.1501 + 1.1145} = \frac{1,446.3874}{131.2646} \approx 11.0188.
	\end{align}
	Taking the conservative integer part so as not to overestimate confidence, we obtain $\nu = 11$ degrees of freedom.
	
	\textbf{Step 2: Critical value from Student's table and standard error.}
	For 99\% confidence ($\alpha = 0.01 \implies \alpha/2 = 0.005$) with $\nu = 11$, the table quantile is $t_{0.005, \, 11} = 3.106$.
	The standard error of the difference between means is:
	\begin{align}
		\text{SE}(\bar{X}_1 - \bar{X}_2) = \sqrt{\frac{S_1^2}{n_1} + \frac{S_2^2}{n_2}} = \sqrt{34.225 + 3.8064} = \sqrt{38.0314} \approx 6.1670 \text{ s.}
	\end{align}
	
	\textbf{Step 3: Construction and interpretation of the 99\% interval:}
	The observed sample difference is $\hat{D} = 120.4 - 105.2 = 15.2$ s. The margin of error $E$ is:
	\begin{align}
		E = 3.106 \times 6.1670 \approx 19.155 \text{ s.}
	\end{align}
	The 99\% interval for $\mu_1 - \mu_2$ is:
	\begin{align}
		15.2 - 19.155 \le \mu_1 - \mu_2 \le 15.2 + 19.155 \implies [-3.955, \, 34.355] \text{ s.}
	\end{align}
	Since it crosses zero from negative to positive ($[-3.96, 34.36]$), we cannot conclude that Algorithm 2 is significantly faster than Algorithm 1 at 99\% confidence.
\end{solproblema}

\begin{solproblema}
	\textbf{Step 1: Degrees of freedom and Fisher ($F$) critical points.}
	Numerator and denominator degrees of freedom are $\nu_1 = n_1 - 1 = 15$ and $\nu_2 = n_2 - 1 = 12$.
	For $98\%$ confidence ($\alpha = 0.02 \implies \alpha/2 = 0.01$), the upper quantile in the $F$ table is:
	\begin{align}
		F_{0.01, \, 15, \, 12} = 3.96.
	\end{align}
	By the geometric reciprocal property for exchanging degrees of freedom, the lower quantile is:
	\begin{align}
		F_{0.99, \, 15, \, 12} = \frac{1}{F_{0.01, \, 12, \, 15}} = \frac{1}{3.67} \approx 0.2725.
	\end{align}
	
	\textbf{Step 2: Calculation of the interval for the population ratio $\dfrac{\sigma_1^2}{\sigma_2^2}$.}
	The empirical ratio between observed sample variances is:
	\begin{align}
		\frac{S_1^2}{S_2^2} = \frac{0.36}{0.12} = 3.00.
	\end{align}
	Substituting into the general formula for the interval of a dispersion ratio:
	\begin{align}
		\frac{S_1^2 / S_2^2}{F_{0.01, \, 15, \, 12}} \le \frac{\sigma_1^2}{\sigma_2^2} \le \frac{S_1^2 / S_2^2}{F_{0.99, \, 15, \, 12}} \implies \frac{3.00}{3.96} \le \frac{\sigma_1^2}{\sigma_2^2} \le \frac{3.00}{0.2725} \implies [0.7576, \, 11.009].
	\end{align}
	\textbf{Interpretation and hypothesis testing conclusion:} The 98\% interval for the ratio of variances spans from $0.758$ to $11.009$. Since the exact equivalence ratio ($1.0$) lies comfortably inside the empirical boundaries, the data demonstrate that there is not enough evidence at the $98\%$ level to conclude that Supplier 1's sensor has higher or lower variability than Supplier 2's.
\end{solproblema}

\subsubsection*{4. Challenging Level (Experts / Deep Deduction)}

\begin{solproblema}
	\begin{enumerate}
		\item \textbf{Theoretical limitations of the symmetric approximation in $r$:}
		The correlation coefficient $r$ is strictly bounded in domain to the compact space $[-1, 1]$. When the true population correlation is strong (for example, $\rho = 0.85$), the sampling distribution of $r$ is strongly compressed against the upper boundary $+1$, generating heavy negative skewness with an elongated left tail. If in this scenario we applied the naive formula $r \pm Z_{\alpha/2}\text{SE}(r)$, the upper limit of the interval would frequently exceed $+1$ (an axiomatic absurdity in correlations), and actual coverage rates would drop drastically below the nominal $95\%$ because the exact variance of the estimator $r$ heavily depends on the parameter $\rho$ itself: $\Var(r) \approx \dfrac{(1-\rho^2)^2}{n}$.
		
		\item \textbf{Inversion of Fisher's hyperbolic transformation:}
		Sir Ronald Fisher derived that applying the inverse hyperbolic tangent to the estimator $r$, the transformed variable $Z(r) = \text{artanh}(r) = \dfrac{1}{2}\ln\left(\dfrac{1+r}{1-r}\right)$ follows an asymptotically normal distribution whose mean is $Z(\rho)$ and, crucially, \textbf{whose variance is stabilized at the constant $\dfrac{1}{n-3}$, independent of the parameter}:
		\begin{align}
			Z(r) \sim N\left( \frac{1}{2}\ln\left(\frac{1+\rho}{1-\rho}\right), \; \frac{1}{n-3} \right).
		\end{align}
		Therefore, the $(1-\alpha)$ symmetric asymptotic interval for parameter $Z(\rho)$ on the transformed scale can be written without error as:
		\begin{align}
			\left[ Z_L, \, Z_U \right] = \left[ Z(r) - Z_{\alpha/2} \frac{1}{\sqrt{n-3}}, \; Z(r) + Z_{\alpha/2} \frac{1}{\sqrt{n-3}} \right].
		\end{align}
		Since the hyperbolic tangent function $\tanh(z) = \dfrac{e^{2z}-1}{e^{2z}+1}$ is strictly monotonically increasing and continuous for all $z \in \mathbb{R}$, it preserves algebraic probability inequalities when applied across all terms:
		\begin{align}
			Z_L \le Z(\rho) \le Z_U \iff \tanh(Z_L) \le \tanh\left(\text{artanh}(\rho)\right) \le \tanh(Z_U) \iff \tanh(Z_L) \le \rho \le \tanh(Z_U).
		\end{align}
		This yields the exact limits on the original scale:
		\begin{align}
			[L_\rho, \, U_\rho] = \left[ \tanh(Z_L), \; \tanh(Z_U) \right].
		\end{align}
		
		\item \textbf{Numerical calculation for the bioinformatics study ($r = 0.82, n = 28$):}
		We evaluate Fisher's transformation at $r = 0.82$:
		\begin{align}
			Z(0.82) = \frac{1}{2} \ln\left( \frac{1 + 0.82}{1 - 0.82} \right) = \frac{1}{2} \ln\left( \frac{1.82}{0.18} \right) = \frac{1}{2} \ln(10.1111) = \frac{1}{2}(2.31363) \approx 1.1568.
		\end{align}
		The constant standard error on the $Z$ scale for $n = 28$ is:
		\begin{align}
			\text{SE}_Z = \frac{1}{\sqrt{n - 3}} = \frac{1}{\sqrt{28 - 3}} = \frac{1}{\sqrt{25}} = \frac{1}{5} = 0.2000.
		\end{align}
		For 95\% confidence ($Z_{0.025} = 1.96$), we compute the intermediate bounds $Z_L$ and $Z_U$:
		\begin{align}
			Z_L &= 1.1568 - 1.96(0.2000) = 1.1568 - 0.3920 = 0.7648, \\
			Z_U &= 1.1568 + 1.96(0.2000) = 1.1568 + 0.3920 = 1.5488.
		\end{align}
		Applying the monotonic inverse function $\tanh(z)$ to each bound to return to the space $[-1, 1]$:
		\begin{align}
			L_\rho &= \tanh(0.7648) = \frac{e^{2(0.7648)} - 1}{e^{2(0.7648)} + 1} = \frac{e^{1.5296} - 1}{e^{1.5296} + 1} = \frac{4.6163 - 1}{4.6163 + 1} = \frac{3.6163}{5.6163} \approx 0.6439. \\
			U_\rho &= \tanh(1.5488) = \frac{e^{2(1.5488)} - 1}{e^{2(1.5488)} + 1} = \frac{e^{3.0976} - 1}{e^{3.0976} + 1} = \frac{22.1448 - 1}{22.1448 + 1} = \frac{21.1448}{23.1448} \approx 0.9136.
		\end{align}
		The exact 95\% asymptotic confidence interval for true population correlation in gene expression is:
		\begin{align}
			IC_{95\%}(\rho) = [0.6439, \; 0.9136].
		\end{align}
		Notice how the resulting interval is logically asymmetric around the point estimate ($0.82 - 0.6439 = 0.1761$ left distance versus $0.9136 - 0.82 = 0.0936$ right distance), perfectly respecting the natural $+1$ boundary.
	\end{enumerate}
\end{solproblema}

\begin{solproblema}
	\begin{enumerate}
		\item \textbf{Algebraic-differential derivation of the Delta Method for ratios:}
		Let the nonlinear objective function of two variables be $g(u, v) = \dfrac{v}{u}$, evaluated at the random vector of sample means $(\bar{X}, \bar{Y})$ whose true expectation parameter is $\boldsymbol{\mu} = (\mu_X, \mu_Y)$.
		
		We expand the first-order bivariate Taylor series of $g(\bar{X}, \bar{Y})$ around the parameter point $(\mu_X, \mu_Y)$:
		\begin{align}
			g(\bar{X}, \bar{Y}) \approx g(\mu_X, \mu_Y) + \left. \frac{\partial g}{\partial u} \right|_{(\mu_X, \mu_Y)} (\bar{X} - \mu_X) + \left. \frac{\partial g}{\partial v} \right|_{(\mu_X, \mu_Y)} (\bar{Y} - \mu_Y).
		\end{align}
		We calculate the exact partial derivatives evaluated at the central point:
		\begin{align}
			\frac{\partial g}{\partial u} &= \frac{\partial}{\partial u}\left( \frac{v}{u} \right) = -\frac{v}{u^2} \implies \left. \frac{\partial g}{\partial u} \right|_{(\mu_X, \mu_Y)} = -\frac{\mu_Y}{\mu_X^2} = -\frac{\theta}{\mu_X}. \\
			\frac{\partial g}{\partial v} &= \frac{\partial}{\partial v}\left( \frac{v}{u} \right) = \frac{1}{u} \implies \left. \frac{\partial g}{\partial v} \right|_{(\mu_X, \mu_Y)} = \frac{1}{\mu_X}.
		\end{align}
		Substituting into the standardized linear expansion of the ratio estimator $\hat{\theta}$:
		\begin{align}
			\hat{\theta} = \frac{\bar{Y}}{\bar{X}} \approx \theta - \frac{\theta}{\mu_X}(\bar{X} - \mu_X) + \frac{1}{\mu_X}(\bar{Y} - \mu_Y).
		\end{align}
		To obtain the asymptotic variance of $\hat{\theta}$, we apply the variance operator to the linear approximation. Recalling that for any linear combination $\Var(aW_1 + bW_2) = a^2\Var(W_1) + b^2\Var(W_2) + 2ab\cov(W_1, W_2)$ and that the constant term $\theta$ has zero variance:
		\begin{align}
			\Var(\hat{\theta}) \approx \left( -\frac{\theta}{\mu_X} \right)^2 \Var(\bar{X}) + \left( \frac{1}{\mu_X} \right)^2 \Var(\bar{Y}) + 2\left( -\frac{\theta}{\mu_X} \right)\left( \frac{1}{\mu_X} \right) \cov(\bar{X}, \bar{Y}).
		\end{align}
		By properties of independent and paired i.i.d. sampling of size $n$, the variances and covariances of the sample averages are:
		\begin{align}
			\Var(\bar{X}) = \frac{\sigma_X^2}{n}, \quad \Var(\bar{Y}) = \frac{\sigma_Y^2}{n}, \quad \cov(\bar{X}, \bar{Y}) = \frac{\sigma_{XY}}{n}.
		\end{align}
		Substituting these sampling variances into the equation and factoring out $1 / (n \mu_X^2)$:
		\begin{align}
			\Var(\hat{\theta}) \approx \frac{\theta^2}{\mu_X^2} \left( \frac{\sigma_X^2}{n} \right) + \frac{1}{\mu_X^2} \left( \frac{\sigma_Y^2}{n} \right) - \frac{2\theta}{\mu_X^2} \left( \frac{\sigma_{XY}}{n} \right) = \frac{1}{n \mu_X^2} \left[ \sigma_Y^2 - 2\theta \sigma_{XY} + \theta^2 \sigma_X^2 \right].
		\end{align}
		This formally proves the fundamental result of the Delta Method for algebraic ratios of averages.
		
		\item \textbf{Systematic interval estimation procedure in practice:}
		Since in empirical data science reality the theoretical parameters $(\mu_X, \sigma_X^2, \sigma_Y^2, \sigma_{XY}, \theta)$ are unknown, to construct the $(1-\alpha)$ confidence interval we follow this sequential protocol:
		\begin{itemize}
			\item Point-estimate the sample averages and unbiased sample variances: $\bar{X}, \bar{Y}, s_X^2, s_Y^2$, and compute the sample covariance of the paired data:
			\begin{align}
				s_{XY} = \frac{1}{n-1} \sum_{i=1}^n (X_i - \bar{X})(Y_i - \bar{Y}).
			\end{align}
			\item Evaluate the central point ratio $\hat{\theta} = \dfrac{\bar{Y}}{\bar{X}}$.
			\item Substitute the estimates into the derived formula to obtain the consistent standard error ($\text{SE}_{\hat{\theta}}$):
			\begin{align}
				\text{SE}(\hat{\theta}) = \sqrt{ \frac{1}{n \bar{X}^2} \left[ s_Y^2 - 2\hat{\theta} s_{XY} + \hat{\theta}^2 s_X^2 \right] }.
			\end{align}
			\item Finally, by asymptotic normality of the linearized estimator as $n \to \infty$, the $(1-\alpha)$ confidence interval for the population ratio $\theta$ is formulated as:
			\begin{align}
				IC_{(1-\alpha)}(\theta) = \hat{\theta} \pm Z_{\alpha/2} \cdot \text{SE}(\hat{\theta}) = \frac{\bar{Y}}{\bar{X}} \pm Z_{\alpha/2} \sqrt{ \frac{1}{n \bar{X}^2} \left[ s_Y^2 - 2\hat{\theta} s_{XY} + \hat{\theta}^2 s_X^2 \right] }.
			\end{align}
		\end{itemize}
	\end{enumerate}
\end{solproblema}
